\documentclass[a4paper,12pt]{amsart}

\usepackage[utf8]{inputenc} % Linux
\usepackage{amsmath}
\usepackage{amsbsy}
\usepackage{amssymb}
\usepackage{amscd,amsthm}
\usepackage{verbatim}
\usepackage[english]{babel}
\usepackage{dsfont}
\usepackage{anysize}
\usepackage{hyperref}

\renewcommand{\proofname}{Proof}
\newcommand{\e}{\varepsilon}
\newcommand{\la}{\lambda}
\newcommand{\tl}{\text{li}}
\newcommand{\Z}{\mathds{Z}}
\newcommand{\Q}{\mathds{Q}}
\newcommand{\N}{\mathds{N}}
\newcommand{\R}{\mathds{R}}
\newcommand{\Pz}{\mathds{P}}
\newcommand{\C}{\mathds{C}}
\newcommand{\lan}{\langle}
\newcommand{\ran}{\rangle}
\newcommand{\ts}{\textsc}
\newcommand{\p}{\phantom}
\newcommand{\q}{\quad}
\newcommand{\al}{\alpha}
\newcommand{\be}{\beta}

\newtheorem{thm}{Theorem}[section]
\newtheorem{lem}[thm]{Lemma}
\newtheorem{kor}[thm]{Corollary}
\newtheorem{conj}[thm]{Conjecture}
\newtheorem{prop}[thm]{Proposition}
\newtheorem*{prim}{Prime $k$-tuples Conjecture}

\theoremstyle{definition}
\newtheorem*{defi}{Definition}
\newtheorem*{nota}{Notation}

\theoremstyle{remark}
\newtheorem*{rema}{Remark}
\newtheorem*{bsp}{Example}

\title{On Ramanujan's prime counting inequality}
\author{Christian Axler}
\address{Institute of Mathematics\\ Heinrich Heine University Duesseldorf\\
40225 Duesseldorf, Germany}
\email{christian.axler@hhu.de}
\date{\today}
\subjclass[2010]{Primary 11N05; Secondary 11M26}
\keywords{Prime counting function, Ramanujan's prime counting inequality, Riemann hypothesis}


\begin{document}

\begin{abstract}
In this paper, we give a new upper bound for the number $N_{\mathcal{R}}$ which is 
defined to be the smallest positive integer such that a certain inequality due to Ramanujan involving the prime counting function $\pi(x)$ holds for every $x \geq N_{\mathcal{R}}$.
\end{abstract}

\maketitle

% \hfill
% \begin{minipage}[t]{50mm}
% \begin{center}
% \textit{Dedicated to the memory of the great mathematician Andrzej Schinzel}
% % \end{minipage}
% \end{center}

\section{Introduction}

Let $\pi(x)$ denote the number of primes not exceeding $x$. Since there are infinitely many primes, we have $\pi(x) \to \infty$ as $x \to \infty$. Gauss \cite{gauss} stated a conjecture concerning an asymptotic behavior of $\pi(x)$, namely
\begin{equation}
\pi(x) \sim \text{li}(x) \quad\quad (x \to \infty), \tag{1.1} \label{1.1}
\end{equation}
where the \emph{logarithmic integral} $\text{li}(x)$ is defined as
\begin{equation}
\text{li}(x) = \int_0^x \frac{\text{d}t}{\log t} = \lim_{\e \to 0} \left \{ \int_{0}^{1-\e}{\frac{\text{d}t}{\log t}} + \int_{1+\e}^{x}{\frac{\text{d}t}{\log 
t}} \right \}, \tag{1.2} \label{1.2}
\end{equation}
where $\log x$ is the natural logarithm of $x$. Hadamard \cite{hadamard1896} and de la Vall\'{e}e-Poussin \cite{vallee1896} independently proved the asymptotic formula \eqref{1.1} which is known as the \textit{Prime Number Theorem}. In a later paper, \cite{vallee1899}
where the existence of a zero-free region for the Riemann zeta-function $\zeta(s)$ to the left of the line $\text{Re}(s) = 1$ was proved, de la Vall\'{e}e-Poussin also estimated the error term in the Prime Number Theorem by showing
\begin{equation}
\pi(x) = \text{li}(x) + O(x \exp(-\delta_0\sqrt{\log x})) \tag{1.3} \label{1.3}
\end{equation} 
as $x \to \infty$, where $\delta_0$ is a positive absolute constant. Integration by parts in \eqref{1.3} implies that for every positive integer $n$, we have
\begin{equation}
\pi(x) = \frac{x}{\log x} + \frac{x}{\log^2 x} + \frac{2x}{\log^3 x} + \frac{6x}{\log^4 x} + \ldots + \frac{(n-1)! x}{\log^nx} + O \left( \frac{x}{\log^{n+1} x} \right) \tag{1.4} \label{1.4}
\end{equation}
as $x \to \infty$. In one of his notebooks (see Berndt \cite{berndt1994}), 
Ramanujan used \eqref{1.4} with $n = 5$ to find that
\begin{displaymath}
\pi(x)^2 - \frac{ex}{\log x} \pi \left( \frac{x}{e} \right)  = - \frac{x^2}{\log^6x} + O \left( \frac{x}{\log^7 x} \right) \tag{1.5} \label{1.5}
\end{displaymath}
as $x \to \infty$ and concluded that the inequality
\begin{equation}
\pi(x)^2 < \frac{ex}{\log x} \pi \left( \frac{x}{e} \right) \tag{1.6} \label{1.6}
\end{equation}
holds for all sufficiently large values of $x$. The inequality \eqref{1.6} is called \textit {Ramanujan's prime counting inequality}.
Recently, Hassani \cite[Corollary 1]{hassani2} found the full asymptotic expansion in \eqref{1.5} by showing that for every integer $n \geq 4$, one has
\begin{displaymath}
\pi(x)^2 - \frac{ex}{\log x} \pi \left( \frac{x}{e} \right) = x^2 \sum_{k=4}^n \frac{d_k}{\log^{k+2}x} + O \left(\frac{x^2}{\log^{n+3} x} \right)
\end{displaymath}
as $x \to \infty$, where
\begin{displaymath}
d_k = \sum_{j=0}^k j! \left( (k-j)! - \binom{k}{j} \right).
\end{displaymath}
A legitimate question is to find the smallest positive integer $N_{\mathcal{R}}$ so that the inequality \eqref{1.6} holds for every real $x \geq N_{\mathcal{R}}$.
The first result made in the search for $N_{\mathcal{R}}$ is based on the assumption that the Riemann hypothesis is true. The Riemann zeta function is for all complex numbers $s$ with $\text{Re}(s) > 1$ defined as
\begin{displaymath}
\zeta(s) = \sum_{n=1}^s \frac{1}{n^s}.
\end{displaymath}
It is well known that the Riemann zeta function is a meromorphic function on the whole complex plane, which is holomorphic everywhere except for a simple pole at $s = 1$.
% Euler discovered a remarkable connection between the Riemann zeta function and the prime numbers by showing that
% \begin{displaymath}
% \zeta(s) = \prod_{p} \frac{1}{1-p^{-s}} \q\q (\text{Re}(s) > 1),
% \end{displaymath}
% where $p$ runs over all prime numbers. Therefore, the Riemann zeta function plays a special role in analytic number theory.
The Riemann zeta function satisfies the functional equation
\begin{displaymath}
\zeta(s)= 2^s \pi^{s-1} \sin \left({\frac{\pi s}{2}} \right) \Gamma(1-s)\zeta (1-s),
\end{displaymath}
where $\Gamma(s)$ is the gamma function. This is an equality of meromorphic functions valid on the whole complex plane. Due to the zeros of the sine function, the functional equation implies that $\zeta(s)$ has outside the set $\{ s \in \C \mid 0 \leq \text{Re}(s) \leq 1 \}$ a simple zero at each even negative integer $s = -2n$, known as the \textit{trivial zeros}. The \textit{nontrivial zeros}, i.e. the zeros in the set $\{ s \in \C \mid 0 \leq \text{Re}(s) \leq 1 \}$, have attracted far more attention because not only is their distribution far less well known, but their study also yields important results concerning primes and related objects in number theory.
% It is known that any non-trivial zero lies in the open strip $\{ s \in \C \mid 0 < \text{Re}(s) < 1 \}$.
The \textit{Riemann hypothesis} asserts that the real part of every nontrivial zero of the Riemann zeta function is $1/2$.
To this day, the Riemann hypothesis is considered one of the greatest unsolved problems in mathematics.
Under the assumption that the Riemann hypothesis is 
true (RH), Hassani \cite[Theorem 1.2]{hassani2012} has given the upper bound 
\begin{displaymath}
RH \; \Rightarrow \; N_{\mathcal{R}} \leq 138\,766\,146\,692\,471\,228.
\end{displaymath}
Dudek and Platt \cite[Lemma 3.2]{dudekplatt} refined Hassani's result by showing that
\begin{equation}
RH \; \Rightarrow \; N_{\mathcal{R}} \leq 1.15 \cdot 10^{16}. \tag{1.7} \label{1.7}
\end{equation}
Wheeler, Keiper and Galway (see Berndt \cite[p.\:113]{berndt1994}) attempted to determine the value of $N_{\mathcal{R}}$, but they failed. Nevertheless, Galway found that 
the largest prime up to $10^{11}$ for which the inequality \eqref{1.6} fails is $x=38\,358\,837\,677$. Hence
\begin{displaymath}
N_{\mathcal{R}} > 38\,358\,837\,677.
\end{displaymath}
Dudek and Platt \cite[Theorem 1.3]{dudekplatt} showed by computation that the largest (not necessarily prime) integer counterexample of \eqref{1.6} up to $x = 10^{11}$ occurs at $x = 38\,358\,837\,682$ and that the inequality \eqref{1.6} holds for every $x$ satisfying $10^{11} \leq x \leq 1.15 \times 10^{16}$. If we combine this result with \eqref{1.7}, it turns out that
\begin{displaymath}
RH \; \Rightarrow \; N_{\mathcal{R}} = 38\,358\,837\,683.
\end{displaymath}
Based on a result of Büthe \cite[Theorem 2]{buethe}, the present author \cite[Theorem 3]{axler} extends the computation of Dudek and Platt by showing that Ramanujan's prime counting inequality \eqref{1.6} holds unconditionally for every $x$ such that $38\,358\,837\,683 \leq x \leq 10^{19}$. This was improved by Platt and Trudgian \cite[Theorem 2]{plattt}. They showed that Ramanujan's prime counting inequality \eqref{1.6} holds unconditionally for every $x$ satisfying $38\,358\,837\,683 \leq x \leq e^{58}$. Recently, Johnston \cite[Theorem 5.1]{johnston} utilized a simple (but computationally intensive) method to show that the inequality \eqref{1.6} holds unconditionally for every $x$ with $38\,358\,837\,683 \leq x \leq e^{103}$.

In another direction, Dudek and Platt \cite[Theorem 1.2]{dudekplatt} claimed to give an upper bound for $N_{\mathcal{R}}$ which does not depend on the assumption that the Riemann hypothesis is true, namely $N_{\mathcal{R}} \leq \exp(9658)$.
%After the present author contacted Platt, one of the authors of \cite{dudekplatt}, to raise some doubts about the
% correctness of the proof of \eqref{1.6}, Platt confirmed after some discussions that the proof of \eqref{1.6} in its
%present form is not correct. This incorrectness motivated the present author to write this paper, where we prove the
%following result. Here, explicit estimates for the prime counting function play an important role.
After the present author identified an error in the proof given in \cite[Theorem 1.2]{dudekplatt}, he proved \cite[Theorem 4]{axler} the even stronger result
\begin{equation}
N_{\mathcal{R}} \leq \exp(9032). \tag{1.8} \label{1.8}
\end{equation}
In the proof of \eqref{1.8}, effective estimates for the prime counting function $\pi(x)$ which hold for all sufficiently large values of $x$ play an important role. Using their effective estimates for $|\vartheta(x)-x|$, where Chebyshev's $\vartheta$-function which is defined by
\begin{equation}
\vartheta(x) = \sum_{p \leq x} \log p, \tag{1.9} \label{1.9}
\end{equation}
where $p$ runs over all primes not exceeding $x$, Platt and Trudgian \cite[Theorem 2]{plattt} fixed the error in the proof given in \cite[Theorem 1.2]{dudekplatt} to show that Ramanujan's prime counting inequality \eqref{1.6} holds unconditionally for every $x \geq \exp(3915)$; i.e. 
\begin{equation}
N_{\mathcal{R}} \leq \exp(3915).  \tag{1.10} \label{1.10}
\end{equation}
Cully-Hugill and Johnston \cite[Corollary 1.6]{cully} used the method investigated by Platt and Trudgian to prove that
\begin{equation}
N_{\mathcal{R}} \leq \exp(3604).  \tag{1.11} \label{1.11}
\end{equation}
Recently, Johnston and Yang \cite[Theorem 1.5]{yang} utilized the same method to improve the last result by showing that
\begin{equation}
N_{\mathcal{R}} \leq \exp(3361).  \tag{1.12} \label{1.12}
\end{equation}
In this paper we will also make use of this method combined with a recent result concerning the difference of $\vartheta(x)$ and $x$ due to Fiori, Kadiri, and Swidinsky \cite{kadiri} to show the following

\begin{thm} \label{thm101}
Ramanujan's prime counting inequality \eqref{1.6} holds unconditionally for every $x \geq \exp(3158.442)$; i.e.
\begin{displaymath}
N_{\mathcal{R}} \leq \exp(3158.442).
\end{displaymath}
\end{thm}

\section{Preliminaries}

The prime counting function $\pi(x)$ and Chebyshev's $\vartheta$-function (cf. \eqref{1.9}) are connected by the 
identity
\begin{equation}
\pi(x) = \frac{\vartheta(x)}{\log x} + \int_{2}^{x}{\frac{\vartheta(t)}{t \log^{2} t}\ \text{d}t}, \tag{2.1} \label{2.1}
\end{equation}
which holds for every $x \geq 2$ (see \cite[Theorem 4.3]{ap}). The method established to prove results like \eqref{1.10}-\eqref{1.12} or Theorem \ref{thm101} goes back to Dudek and Platt \cite{dudekplatt} and is as follows. We start with a function $a$ so that
\begin{equation}
|\vartheta(x)-x| \leq \frac{a(x)x}{\log^5x} \tag{2.2} \label{2.2}
\end{equation}
for every $x \geq x_0$. Then we substitute this inequality into \eqref{2.1} to derive upper and lower bounds for the prime counting function $\pi(x)$ of the form
\begin{equation}
x\sum_{k=0}^4 \frac{k!}{\log^{k+1}x} + \frac{m_a(x)x}{\log^6x} \leq \pi(x) \leq x\sum_{k=0}^4 \frac{k!}{\log^{k+1}x} + \frac{M_a(x)x}{\log^6x} \tag{2.3} \label{2.3}
\end{equation}
for every $x \geq x_1$, where the functions $m_a(x)$ and $M_a(x)$. Using these estimates, Dudek and Platt were able to give an
explicit version of \eqref{1.5}. As already mentioned in the introduction, we also use this method. So, we need to find a function $a$ and a positive real number $x_0$ so that the inequality \eqref{2.2} holds. For this purpose, we set
\begin{displaymath}
R = 5.5666305
\end{displaymath}
and, similar to \cite[p.\:879]{plattt}, we define the function $a: \mathbb{R}_{>0} \to \mathbb{R}$ by
\begin{displaymath}
\frac{a(x)}{\log^5x} = 
\begin{cases}
\displaystyle \frac{2 - \log2}{2} & \text{if $2 \leq x < 599$,} \\
\displaystyle\frac{\log^2x}{8\pi\sqrt{x}} & \text{if $599 \leq x < 1.101 \times 10^{26}$,} \\
\displaystyle \sqrt{\frac{8}{17\pi}} \left( \frac{\log x}{6.455} \right)^{1/4} \exp \left( - \sqrt{\frac{\log x}{6.455}} \right) & \text{if $1.101 \times 10^{26} \leq x < e^{673}$,} \\
\displaystyle 121.0961 \left( \frac{\log x}{R} \right)^{3/2} \exp \left( - 2 \sqrt{\frac{\log x}{R}} \right) & \text{if $x \geq e^{673}$.}
\end{cases}
\end{displaymath}
Then we get the following result concerning Chebyshev's $\vartheta$-function.

\begin{lem} \label{lem201}
For every $x \geq 2$, we have
\begin{displaymath}
|\vartheta(x) - x| \leq \frac{a(x)x}{\log^5x}.
\end{displaymath}
\end{lem}

\begin{proof}
If $x$ satisfies $2 \leq x < 599$, then the given bound is trivial. The second one was proven by Johnston \cite[Corollary 3.3]{johnston} and the third bound was given by Trudgian \cite[Theorem 1]{trud}. The last bound was recently established by Fiori, Kadiri, and Swidinsky \cite[Corollary 14]{kadiri}.
\end{proof}

We also need the following result on our function $a$.

% Next we note that the function $a_n(x)$ is a decreasing function in the has a global maximum at $x_0 = e^{(4n+1)^2R/4}$.

\begin{lem} \label{lem202}
Let $x_1$ be real number with $x_1 \geq e^{673}$. Then $a_n(x) \leq a_n(x_1)$ for every $x \geq x_1$.
\end{lem}

\begin{proof}
By a straightforward calculation of the derivative, we see that the inequality $a'(x) < 0$ holds for every $x \geq e^{673}$.
\end{proof}

Now, let $x_1$ be a real number with $x_1 \geq e^{673}$. According to the method we use, we set
\begin{align*}
C_0 &= \int_2^{x_1} \frac{720 - a(t)}{\log^7 t} \, \text{d}t - 2\sum_{k=1}^5 \frac{k!}{\log^{k+1}2}, \\
C_1 &= \int_2^{x_1} \frac{720 + a(t)}{\log^7 t} \, \text{d}t - 2\sum_{k=1}^5 \frac{k!}{\log^{k+1}2}. 
\end{align*}
Contrary to Dudek and Platt \cite{dudekplatt}, who have estimated the integral
\begin{displaymath}
\int_2^x \frac{\text{d}t}{\log^7t},
\end{displaymath}
we will use the identity
\begin{equation}
\int_{x_1}^x \frac{\text{d}t}{\log^7t} = E(x) - E(x_1), \tag{2.4} \label{2.4}
\end{equation}
where
\begin{displaymath}
E(x) = \frac{1}{720} \left( \text{li}(x) - \frac{x}{\log x} - \frac{x}{\log^2x} - \frac{2x}{\log^3x} - \frac{6x}{\log^4x} - \frac{24x}{\log^5x} - \frac{120x}{\log^6x}\right),
\end{displaymath}
to find the following explicit version of \eqref{2.3}.

\begin{lem} \label{lem203}
Let
\begin{displaymath}
m_a(x) = 120 - a(x) + \left( C_0 + (720 - a(x_1))(E(x) - E(x_1)) \right) \frac{\log^6x}{x}
\end{displaymath}
and 
\begin{displaymath}
M_a(x) = 120 + a(x) + \left( C_1 + (720 + a(x_1))(E(x) - E(x_1)) \right) \frac{\log^6x}{x}.
\end{displaymath}
Then
\begin{displaymath}
x\sum_{k=0}^4 \frac{k!}{\log^{k+1}x} + \frac{m_a(x)x}{\log^6x} \leq \pi(x) \leq x\sum_{k=0}^4 \frac{k!}{\log^{k+1}x} + \frac{M_a(x)x}{\log^6x}
\end{displaymath}
for every $x \geq x_1$.
\end{lem}

\begin{proof}
We only give a proof of the required upper bound. The proof of the required lower bound is similar and we leave the details to the reader. Let $x \geq x_1$. If we combine \eqref{2.1} with Lemma \ref{lem201}, we can see that
\begin{displaymath}
\pi(x) \leq \frac{x}{\log x} + \frac{xa(x)}{\log^6x} + \int_2^x \frac{\text{d}t}{\log^2 t} + \int_2^x \frac{a(t)}{\log^7 t} \, \text{d}t.
\end{displaymath}
Integration by parts in \eqref{1.3} provides that
\begin{displaymath}
\pi(x) \leq x\sum_{k=0}^4 \frac{k!}{\log^{k+1}x} + \frac{(120+a(x))x}{\log^6x} + C_1 + \int_{x_1}^x \frac{720 + a(t)}{\log^7 t} \, \text{d}t.
\end{displaymath}
Since $a(t) \leq a(x_1)$ for every $t$ with $x_1 \leq t \leq x$ (cf. Lemma \ref{lem202}), it turns out that
\begin{displaymath}
\pi(x) \leq x\sum_{k=0}^4 \frac{k!}{\log^{k+1}x} + \frac{(120+a(x))x}{\log^6x} + C_1 + (720 + a(x_1))\int_{x_1}^x \frac{\text{d}t}{\log^7 t}.
\end{displaymath}
Finally, it suffices to apply the identity \eqref{2.4}.
\end{proof}

\section{Proof of Theorem \ref{thm101}}

Now we have all the necessary tools to give a proof of Theorem \ref{thm101}.

\begin{proof}[Proof of Theorem \ref{thm101}]
Let $x_1 = \exp(3157.442)$. Then, one has $a(x_1) = 1056.767676\ldots$ and $E(x_1) > 0$. Since the function $E(x)\log^6(x)/x$ is decreasing on the interval $[x_1, \infty)$ and $a(x) \leq a(x_1)$ for every $x \geq x_1$ (cf. Lemma \ref{lem202}), we can use Lemma \ref{lem203} and a computer to get that
\begin{displaymath}
x\sum_{k=0}^4 \frac{k!}{\log^{k+1}x} + \frac{m_a x}{\log^6x} \leq \pi(x) \leq x\sum_{k=0}^4 \frac{k!}{\log^{k+1}x} + \frac{M_a x}{\log^6x}
\end{displaymath}
for every $x \geq x_1$, where
\begin{align*}
m_a &= -936.64603213534,\\
M_a &= 1177.56019022252.
\end{align*}
Now we can argue as in the proof of \cite[Lemma 2.1]{dudekplatt} to see that
\begin{equation}
\pi^2(x) - \frac{ex}{\log x}\, \pi \left( \frac{x}{e} \right) < \frac{x^2}{\log^6x} \left( -1 + \frac{\epsilon_{M_a}(x) - \epsilon_{m_a}(x)}{\log x} \right) \tag{3.1} \label{3.1}
\end{equation}
for every $x \geq ex_1$, where
\begin{align*}
\epsilon_{m_a}(x) &= 206 + m_a + \frac{364}{\log x} + \frac{381}{\log^2x} + \frac{238}{\log^3x} + \frac{97}{\log^4x} + \frac{30}{\log^5x} + \frac{8}{\log^6x}, \\
\epsilon_{M_a}(x) &= 72 + 2M_a + \frac{2M_a + 132}{\log x} + \frac{4M_a + 288}{\log^2 x} + \frac{12M_a + 576}{\log^3 x} + + \frac{48M_a}{\log^4 x} + \frac{M_a^2}{\log^5x}.
\end{align*}
Note that $\epsilon_{M_a}(x) - \epsilon_{m_a}(x) < \log x$ for every $x \geq ex_1 = \exp(3158.442)$. Finally, it suffices to substitute the last inequality into \eqref{3.1} and we arrive at the end of the proof.
\end{proof}

\section{Future work}
It is natural to ask whether we can derive comparable results if we replace the number $e$ in \eqref{1.6} by an arbitrary positive real number $\alpha$. In this context, Hassani \cite[Theorem 3]{hassani2} was able to show that if $\alpha > e$ then one has
\begin{displaymath}
\pi(x)^2 < \frac{\alpha x}{\log x} \pi \left( \frac{x}{\alpha} \right) \tag{4.1} \label{4.1}
\end{displaymath}
for all sufficiently large values of $x$ and if $0 < \alpha < e$, then the above inequality reverses. One could inverstigated a method, similar to the one we used in the proof of Theorem \ref{thm101}, to find effective estimates for the smallest positive integer $N_{\mathcal{R}, \alpha}$ so that the inequality \eqref{4.1} holds for every real $x \geq N_{\mathcal{R}, \alpha}$.

\section*{Acknowledgement}

The author wishes to thank the two beautiful souls R. and O. for the never ending inspiration.
% Futhermore, the author would like to express his great appreciation to the anonymous referees for useful comments to improve the quality of this paper.

\begin{thebibliography}{10}
\bibitem{ap} T. Apostol, {\it Introduction to Analytic Number Theory}, Springer, New York--Heidelberg, 1976.
\bibitem{axler} C. Axler, \textit{Estimates for $\pi(x)$ for large values of $x$ and Ramanujan's prime counting inequality}, Integers \textbf{18} (2018), Paper No. A61, 14 pp.
\bibitem{berndt1994} B. C. Berndt, {\it Ramanujan's Notebooks, Part IV}, Springer, New York, 1994.
\bibitem{buethe} J. Büthe, \textit{An analytic method for bounding $\psi(x)$}, Math. Comp. \textbf{87} (2018), no. 312, 1991--2009.
\bibitem{cully} M. Cully-Hugill and D. R. Johnston, \textit{On the error term in the explicit formula of Riemann--von Mangoldt}, preprint, 2021. Available at \url{arxiv.org/abs/2111.10001}.
\bibitem{vallee1896} C.-J. de la Vallée Poussin, \textit{Recherches analytiques la th\'{e}orie des nombres premiers}, Ann. Soc. scient. Bruxelles \textbf{20} (1896), 183--256.
\bibitem{vallee1899} C.-J. de la Vallée Poussin, \textit{Sur la fonction $\zeta(s)$ de Riemann et le nombre des nombres premiers inf\'{e}rieurs \`{a} une limite donn\'{e}e}, Mem. Couronn\'{e}s de l'Acad. Roy. Sci. Bruxelles \textbf{59} (1899), 1--74.
\bibitem{dudekplatt} A. W. Dudek and D. J. Platt, \textit{On solving a curious inequality of Ramanujan}, Exp. Math. {\bf 24} (2015), no. 3, 289-294.
\bibitem{kadiri} A. Fiori, H. Kadiri, and J. Swindisky, \textit{Sharper bounds for the error term in the Prime Number Theorem}, preprint, 2022. Available at \url{arxiv.org/abs/2206.12557}.
\bibitem{gauss} C. F. Gauss, \textit{Werke}, 2 ed., Königlichen Gesellschaft der Wissenschaften, Göttingen, 1876.
\bibitem{hadamard1896} J. Hadamard, \textit{Sur la distribution des z\'{e}ros de la fonction $\zeta(s)$ et ses cons\'{e}quences arithm\'{e}tiques}, Bull.Soc. Math. France \textbf{24} (1896), 199--220. 
\bibitem{hassani2012} M. Hassani, \textit{On an inequality of Ramanujan concerning the prime counting function}, Ramanujan J. {\bf 28} (2012), no. 3, 435-442.
\bibitem{hassani2} M. Hassani, \textit{Remarks on Ramanujan's inequality concerning the prime counting function}, Commun. Math. \textbf{29} (2021), no. 3, 473--482.
\bibitem{johnston} D. R. Johnston, \textit{Improving bounds on prime counting functions by partial verification of the Riemann hypothesis}, preprint, 2021. Available at \url{arxiv.org/abs/2109.02249}.
\bibitem{yang} D. R. Johnston and A. Yang, \textit{Some explicit estimates for the error term in the prime number theorem}, preprint, 2022. Available at \url{arxiv.org/abs/2204.01980}.
\bibitem{plattt} D. J. Platt and T. S. Trudgian, \textit{The error term in the prime number theorem}, Math. Comp. \textbf{90} (2021), no. 328, 871-–881.
\bibitem{trud} T. S. Trudgian, \textit{Updating the error term in the prime number theorem}, Ramanujan J. {\bf 39} (2016), no. 2, 225-234.
\end{thebibliography}


\end{document}
